
\chapter*{General Introduction}
\markboth{General Introduction}{} %pour afficher l'entete
\addcontentsline{toc}{chapter}{General introduction}

In an era marked by rapid technological advancement, artificial intelligence, and exponential data growth, organizations face increasing challenges in transforming information into actionable knowledge. AI-powered systems are reshaping how professionals interact with complex data, enabling faster analysis and more informed decision-making.

Within this evolving landscape, Jade Advisory, a consulting firm specializing in Public-Private Partnerships (PPP), manages a large and diverse portfolio of strategic documents. These documents are essential for project execution and compliance. However, their volume and heterogeneity make information retrieval slow, fragmented, and difficult to standardize.

Recognizing the value of its institutional knowledge, Jade Advisory identified the need for an intelligent solution to structure and efficiently access its document repository. The lack of such a system limits knowledge reuse and reduces operational efficiency.

To address this challenge, this project proposes the development of an intelligent document management system based on a Retrieval-Augmented Generation (RAG) approach. Relying on cutting-edge technologies in artificial intelligence and data engineering, the proposed solution transforms the document repository into a structured and searchable knowledge base, enabling accurate information retrieval and question-answering. By leveraging advanced AI techniques, the system provides fast and context-aware access to organizational knowledge, thereby enhancing productivity and supporting informed decision-making.

This report has been prepared as part of the end-of-study internship and in fulfillment of the requirements for the degree in Software Engineering. It is organized into five chapters, reflecting the methodological approach followed during its development :
\begin{itemize}
    \item \textbf{Chapter 1 titled "Project Context" }first presents the host company, then presents the existing, points out its limitations, and finally presents a proposed solution.
    \item \textbf{Chapter 2 titled "Needs Analysis"} presents the functional needs of the proposed solution and the technical specification including non-functionnal need, Hardware environment, Software environment, used technologies and used libraries.
\end{itemize}



